\documentclass[11pt,a4paper]{article}
\usepackage[margin=26mm,headheight=15pt]{geometry}
\usepackage{fontspec}\setmainfont{DejaVu Serif}\setsansfont{DejaVu Sans}
\usepackage{amsmath,amssymb,amsthm,booktabs,longtable,array,microtype,xcolor,fancyhdr,pdfpages}
\usepackage{unicode-math}\setmathfont{latinmodern-math.otf}
\usepackage[colorlinks=true,linkcolor=teal!55!black,urlcolor=teal!55!black,citecolor=teal!55!black]{hyperref}
\definecolor{ink}{HTML}{173C4C}
\pagestyle{fancy}\fancyhf{}\fancyhead[L]{\small\sffamily Scientific assessment and further results}\fancyhead[R]{\small Collections}\fancyfoot[C]{\thepage}
\setlength{\emergencystretch}{2em}
\setlength{\parindent}{0pt}\setlength{\parskip}{5pt}
\title{\color{ink}\sffamily From Prephysical Hypotheses to Physical Models\\[7pt]\Large Collections\\[7pt]\large Scientific assessment and further results}
\author{Kirtus Justus\\\small Port Coquitlam, British Columbia, Canada}
\date{10 November 2026}
\begin{document}
\maketitle
\begin{abstract}
This assessment examines the supplied 81-page Collections and 420-page Mathematical Supplement, their source package and their account of verification. Six distributed computational procedures reproduce their stored results exactly. The 40 original-manuscript titles and 215 source-edition article pages are identified, as are all 80 current components and 33 visible supporting-proof items. Selected entropy, composition and abstract-channel statements withstand the concentrated checks. Two historical Wilson scope sentences remain inconsistent with later included results. The explicit Wilson estimates also yield a conditional uniform additive remainder on bounded scaled-mass intervals, while leaving relative vanishing-mass asymptotics separate. A companion mathematical study gives exact joint treatment of finite calibration and source samples, a concrete undercoverage example, and a limit on uniform inference near zero readout gain. These findings support specified mathematical statements and reproducibility claims. They supply no empirical identification of a prephysical substrate, no independent peer review, and no certification of the complete historical proof corpus.
\end{abstract}
\tableofcontents
\newpage
\section{Scope, source basis and evidential standard}

The materials examined are the user-supplied review packet, its two compiled volumes, the one-page figure \emph{Structure, observation and physical identification}, and the preserved source editions needed for targeted comparisons. The packet has 341 files. The two main volumes contain 501 pages. The assessment distinguishes actual source inspection, code execution, exact finite arithmetic, analytic proof dependencies, editorial consistency and empirical evidence.

The computational and inventory findings below concern the exact supplied versions. Full file identities, procedure hashes, commands, output comparisons and a page concordance accompany this assessment. A successful reproduction checks the named computation. It does not certify every proposition cited by the program, the physical interpretation of the inputs, or an earlier execution by another researcher.

The mathematical review concentrated on resources and equality patterns, relational observations, and the fixed-positive-scaled-mass Wilson chain. It did not freshly verify all 501 pages or all archived programs. No inference about an uninspected proof is made from the number of files, pages or predicates. The review is AI-assisted work within the same research program, with shared source and model lineage; it is not independent external review.

\section{Claims that the new checks support}

\subsection{Diagram scope}
The supplied figure no longer presents a candidate channel construction as the ontology of the entire program. It separates proposed relational structures, an open physical correspondence, specified observations, and conditional mathematical implications. The dashed arrow is explicitly an unresolved correspondence. The lower branches describe possible conclusions of future tests, with calibration assumptions retained. Its references name manuscripts rather than an image filename.

The figure is reproduced at the end of this volume. Its five examples illustrate observation problems; they neither exhaust the program nor establish a unique ontology. Checking this scope statement does not independently certify every formula summarized in every box.

\subsection{Computational reproduction}
Each of six named procedures was inspected and run from an isolated unchanged copy, without giving it the stored comparison result as input. The readout verifier received its specified rational certificate. The environment matched the recorded Python 3.12.14, SymPy 1.14.0 and mpmath 1.3.0 versions. All six completed successfully and reproduced their comparison JSON byte for byte; the submitted files remained unchanged.

\begin{center}\small
\begin{tabular}{p{84mm}r}\toprule
Procedure & Reported predicates\\\midrule
Graph moment certificate & Separate certificate\\
Thirty-six-mass readout transport & Separate certificate\\
Mathematical examples & 680\\
Operational results & 418\\
Finite signs & 449\\
Coverage additions & 1,735\\\midrule
Four-suite total & 3,282\\\bottomrule
\end{tabular}
\end{center}

The count is accurate, but its units vary. One predicate may compare a scalar, a matrix, a probability law, or an entire collection of cases. For example, 1,536 predicates concern one conditional-odds identity over the specified four-vertex graphs, vertex pairs and outside assignments. In the operational suite, 192 predicates are normalization checks and 182 exclude smaller horizons for a fixed testing problem. The count is not a count of independent theorems or proofs. The supplied scope note already states this limitation.

The graph verifier checks a rational contraction construction and a coarse fourth-moment enclosure. Its characteristic-polynomial brackets concern the rational centre; transfer to the nearby root and the lifted graph uses the manuscript's perturbation and invariant-subspace arguments. The transport verifier checks the 36-entry joint probability law, all stated marginals and conditional flip constraints by exact rational arithmetic. It does not thereby establish nonanticipation, a physical causal mechanism, or a history-conditional error bound.

The two exponential tail evaluations for the 8,982-record example reproduce the stated values at 75-digit working precision. They are ordinary multiprecision evaluations, not outward-rounded interval certificates. The exact rational transport proof and those evaluations remain separate claims. No stricter numerical-certification status is assigned by this review.

\subsection{File and article coverage}
All 340 file-manifest entries match their declared sizes and digests; the manifest itself is the documented unlisted file. All 80 component PDF digests and page counts match, and their declared ranges are contiguous within the 81- and 420-page volumes.

The source-edition count of 215 original article pages is correct for the preserved editions: six articles occupy pages 28–65 of the 65-page Collections (38 pages), and 34 articles occupy pages 46–222 of the 234-page Mathematical Review (177 pages). All forty titles match the corresponding current slots, allowing punctuation normalization and the explicit collected-article subtitle. Those slots occupy 226 pages in the new typesetting. The difference is an edition distinction, not evidence of omitted articles.

The 33 supporting-proof slots occupy 104 visible pages; six additional article slots occupy 77 pages; the terminology unit occupies 11 pages. Together with 79 study pages and four contents pages, these account for the 501 pages. This inventory establishes inclusion and pagination. It does not establish that every excerpt contains a complete proof or that all 36 claimed subject areas have been independently assessed for semantic coverage.

The exact original source versions of the graph and causal-readout programs were recovered. Their hashes match the original-source identities declared in the computational note. Their abstract syntax agrees with the distributed programs after documentation strings and comments are excluded. The present run reproduces the distributed programs. Historical execution events and the broader claimed comparison of 49 programs have not been reconstructed here.

The portable assembly procedure was inspected, but no second complete assembly run was performed in this assessment. It assembles supplied compiled components rather than recompiling every historical manuscript. The other preparation's assembly result must not be counted as an additional run by this review.

\section{Mathematical findings that survive concentrated scrutiny}

\subsection{Abstract information does not require physical duration}
The general objection that a channel or Gaussian information expression must already presuppose physical time, power and apparatus is too strong. Abstract channel theory specifies input and output spaces, probability laws and cost constraints \cite{polyanskiy}. For example,
\[
Y=X+Z,\quad Z\sim\mathcal N(0,\sigma^2I_d),\quad Z\perp X,\quad
\mathbb E\|X\|^2\le dP
\]
has information supremum
\[
\sup I(X;Y)=\frac d2\log_2(1+P/\sigma^2)
\]
in bits per real vector block. The entropy bound and Gaussian attainment are mathematical statements. Coding conclusions require the specified repeated-use experiment. An engineering bandwidth and rate in bits per second require further physical assignments.

The new packet permits mathematical indices and spectral variables without identifying them as physical frequency. That correction is sound. It leaves a different obligation: the proposed prephysical assumptions must either derive or openly supply the probability law, input class, noise, cost and observation map. Relabelling those ingredients does not derive them. An unspecified signal-to-noise variance ratio alone does not imply Gaussian capacity; nor does a mathematical representation prove its physical correspondence.

\subsection{Entropy finiteness and attainment}
The resources study correctly states an entropy \emph{supremum}. On its finite/countable alphabet, with finite nonnegative costs, a zero-cost symbol, all input laws and a positive finite mean-cost budget, finiteness of the supremum is equivalent to convergence of a Gibbs partition sum at some positive inverse temperature. Attainment is an additional question.

A direct adverse control illustrates the distinction. Take one zero-cost symbol and $g_n=\lfloor2^n/n^3\rfloor$ symbols of cost $n$. At inverse temperature $\log2$, the Gibbs partition sum is finite and the Gibbs mean cost is strictly below one. At budget one, distributions mixing that Gibbs law with a small mass on increasingly high cost levels approach the finite entropy bound without attaining it. Thus replacing “supremum” by “maximum” would be false in the generic theorem. This construction supports the correction; it does not refute the corrected statement. The role of stronger Gibbs hypotheses in entropy continuity and attainment is also explicit in current mathematical literature \cite{becker}.

The integer-partition example has the additional structure needed for a maximizer. Its partition product converges for every positive inverse temperature, and its Gibbs mean decreases continuously from infinity to zero. The asserted mean-one Gibbs optimum therefore survives this objection. The focused checks found no decisive false theorem in these resource statements.

\subsection{Composition and hidden coupling}
The composition results retain the essential distinction between fresh conditional kernel sampling and persistent shared randomness. The supplied finite nonassociativity witnesses reproduce exactly. A persistent latent random choice can change the composite experiment while preserving the one-call law, so it cannot be substituted silently into a theorem about fresh kernel composition.

Two minor qualifications would improve the written scope. The deterministic natural-operation theorem should say that the operation arity is a \emph{finite integer}. An infinite-arity cofinite-equality operation supplies a different class. The common-reference information identity written as a sum should specify a discrete/countable message law, or be rewritten as a conditional integral with the appropriate measurable-space assumptions. These are scope clarifications under the apparent intended conventions, not identified counterexamples to the finite theorems.

\section{Remaining Wilson inconsistencies and a stronger conditional implication}

The fixed-positive-scaled-mass finite-part paper and the uniform modular paper are not new proof texts in this submission: their eight and five extracted-text pages respectively match the copies already present in the preserved mathematical materials. The substantive change is how the synthesis reports their status.

Two older included passages still contradict that later status:

\begin{enumerate}
\item \emph{Exact Spectral Majorants for a Wilson–Dirac Interval}, component p. 5, current Mathematical Supplement p. 176, says that the full constant-term expansion remains conditional. Its own majorant theorem may properly depend on separate convergence estimates, but the sentence should now identify the later finite-part and modular papers that supply the claimed estimates for fixed $\mu>0$.
\item \emph{Fourth-Moment Control of the Wilson Nonlinear Correction}, component p. 3, current Mathematical Supplement p. 190, says that a full-spectrum $\|E_N\|=O(1)$ bound remains open. The later \emph{Uniform Modular Control and Nonlinear Convergence} asserts $\|E_N\|\le7$ for every even $N\ge2$, using its exact concentration connection.
\end{enumerate}

These are concrete editorial inconsistencies. They do not refute the local majorant or fourth-moment arguments. Suitable replacement wording is:

\begin{quote}
The present paper supplies the majorant, pairing and spectral-tail inputs. The later finite-part and uniform modular manuscripts provide the additional convergence estimates for the specified infinite-chain model at fixed positive scaled mass; their analytic dependencies remain those stated in those papers.
\end{quote}

\begin{quote}
The present reduction uses only the earlier modular bound. The later uniform modular manuscript proves a full-spectrum bound of seven for even dimensions under its stated exact concentration connection; sharper asymptotics and other model regimes are separate questions.
\end{quote}

There is also a useful quantifier correction. The fixed-$\mu$ notation alone does not establish uniformity as $\mu\to0$. However, the displayed finite estimates contain explicit mass dependence. Combined with the already stated uniform boundary estimate and the exact two-bond formula, they imply a \emph{uniform additive} expansion for $0\le\mu\le M$. The companion supplement supplies the full conditional derivation:
\begin{align*}
\sup_{0\le\mu\le M}\left|\mathcal D_L(\mu)
-\frac{\mu^2}{6}\log L
-\mu^2\left(\frac{13}{36}-\frac\pi{12}-\frac{\log\mu}{6}\right)
+\mathcal S(\mu)\right|
=O_M(1/\log L).
\end{align*}
All products at zero are interpreted continuously. The bound is conditional on the supplied exact spectral, majorant and boundary estimates; this review has not independently certified every operator-domain and regularity input.

The four contribution rates remain separate: modular $O_M(\sqrt{\log N}/N)$, nonlinear $O_M(\sqrt{\log N/N})$, two-bond $O_M(\log L/L)$ and boundary $O_M(1/\log L)$. The boundary term still controls this overall additive remainder. No cancellation is assumed. The two square-root expressions are different rates.

This additive result applies along bounded vanishing-mass sequences, but it does not prove $\mathcal D_L(\mu_L)\sim\mu_L^2\log L/6$ for arbitrary $\mu_L\to0$. The relative error need not vanish, and the finite part itself requires relative control. The claim that no uniform statement can follow is therefore too broad; a general relative leading-law assertion is still unsupported by these estimates.

\section{Joint calibration and sampling: the further result}

The companion supplement formalizes the population identified-set map
\[
\mathcal I(P)=\{h(\xi):\mathcal O(\xi)=P,\ \xi\text{ belongs to the admitted model class}\}.
\]
If a random region contains the full true observation law with probability at least $1-\alpha$, the union of $\mathcal I(Q)$ over that region contains the entire true identified set on the same event. This is a sufficient construction. It may be conservative, and it does not turn an unjustified calibration range into a confidence statement. Whole-set and point-parameter coverage remain different targets \cite{romano}.

For the binary readout $p=(1-\theta)f+\theta s$, the exact projection of a rectangular law region is
\[
0\le\theta\le1,\qquad
(1-\theta)f_-+\theta s_-\le p_+,\qquad
(1-\theta)f_++\theta s_+\ge p_-.
\]
It is an interval or the empty set, with zero slopes handled directly. Exact $p=1/2$ and the calibration ranges $f\in[1/5,3/10]$, $s\in[7/10,4/5]$ leave the population interval $[2/5,3/5]$. More source samples cannot remove that calibration ambiguity by themselves.

\subsection{Exact constructed experiment}
The finite experiment uses 80 source trials and four trials for each of two reference preparations. Its fixed parameter family has $\theta=k/6$, $k=0,\ldots,6$, and $f,s\in\{1/5,1/2,4/5\}$. All three binomial streams are stipulated independent. The 63 parameter points induce 45 distinct observation laws. The procedures, grid and sample sizes were fixed before execution and were not tuned to the result.

The joint procedure assigns one third of a 5\% error budget to each stream and projects the admitted finite parameter triples. The plug-in procedure treats the two observed reference fractions as exact and accounts only for source uncertainty. A known-calibration control uses the true reference probabilities. Every count outcome is enumerated with exact rational probabilities.

\begin{center}\small
\begin{tabular}{p{73mm}rr}\toprule
Procedure & Minimum set coverage & Laws below 95\%\\\midrule
Joint calibration and source uncertainty & 98.6328\% & 0/45\\
Reference estimates treated as exact & 13.6591\% & 45/45\\
Known calibration & 95.4497\% & 0/45\\\bottomrule
\end{tabular}
\end{center}

The plug-in minimum whole-set value occurs at an uninformative readout. The failure also occurs in an interior, identified example: for $\theta=1/2$, $f=1/5$, $s=4/5$, the plug-in coverage is 56.2424\%, versus 96.7007\% for known calibration. Joint coverage is 99.9994\%, but its expected confidence-set size is 6.9821 of the seven grid values. This example therefore demonstrates both the danger of omitted calibration uncertainty and the cost of honest uncertainty with very small reference samples. It does not establish that the joint construction is optimal or useful for precision at these sample sizes.

A second implementation, using rational probability recurrences and conditional count-interval calculations rather than the primary integer-weight masks, agrees on all 135 whole-set coverages, 189 admissible point coverages and 45 true-law admission probabilities. These 369 comparisons verify the declared arithmetic. The two implementations share the model, result exposure and program lineage; their agreement is not an independent scientific study.

\subsection{A general limitation near zero gain}
For $f=1/2-\delta$, $s=1/2+\delta$, the two hypotheses $\theta=0$ and $\theta=1$ have identical reference laws. Their $N$ source trials differ in total variation by at most $2N\delta$. Any confidence set with coverage at least $1-\alpha$ for both hypotheses must, under the first law, contain both endpoints with probability at least $1-2\alpha-2N\delta$.

Taking $\delta_N=N^{-2}$ shows why point identification at every fixed nonzero gain does not imply uniformly precise inference over arbitrarily small gains. At 95\% confidence the lower bound on reporting diameter one approaches 90\%. This is a proof about the stated family, independent of the finite grid. It does not rule out consistent inference when a warranted positive gain bound is available.

\section{Consequences for the physical research question}

The checked material improves the precision of mathematical questions and exposes where observation loses information. It does not yet warrant a unique physical ontology. A proposed relational structure, a signed probability law, a quantum completion and a dynamics can be distinct mathematical constructions, each with its own assumptions and observable consequences. Their order in a document is not a derivation connecting them.

For each candidate, a substantive next result would specify a target invariant, an admissible preparation and observation class, a calibration-data model, and an observation that separates the candidate from a concrete competitor. It must then distinguish population identification from finite-data discrimination, and fixed calibration from estimated or drifting calibration. The joint projection principle applies after those ingredients are supplied; the binary numbers in this assessment do not transfer automatically to a quantum instrument or a controlled realization.

Equally, a proof gap is not a counterexample, an observation-equivalent description is not automatically a different physical world, and a scoped conditional theorem remains a valid mathematical achievement if its implication is correct. The current assessment supports several such implications, identifies two remaining editorial corrections, and supplies quantitative limits on inference. It has found no new empirical result that discriminates the proposed prephysical substrate from its competitors.

\section*{Preparation and use}
This volume and its companion are additions to the supplied research materials. They preserve the source editions and identify the exact extent of checking. They do not replace the full supplied proof corpus or represent an assertion that it has passed external review. AI assistance was used for derivation, source comparison, computation and preparation; all computational and analytic limitations stated above remain part of the findings.

\begin{thebibliography}{9}
\bibitem{polyanskiy} Y. Polyanskiy and Y. Wu, \emph{Information Theory}, MIT 6.441 (2016), Chapter 17, “Channels with input constraints. Gaussian channels.” Definitions 17.1–17.5 and Theorem 17.5. \url{https://ocw.mit.edu/courses/6-441-information-theory-spring-2016/}.
\bibitem{becker} S. Becker, N. Datta, M. G. Jabbour and M. E. Shirokov, \emph{Optimal continuity bound for the von Neumann entropy under energy constraints}, arXiv:2410.02686v2 (2025), section 2. \url{https://arxiv.org/abs/2410.02686}.
\bibitem{romano} J. P. Romano and A. M. Shaikh, “Inference for the identified set in partially identified econometric models,” \emph{Econometrica} 78(1), 169–211 (2010). \href{https://doi.org/10.3982/ECTA6706}{doi:10.3982/ECTA6706}.
\bibitem{supplied} K. Justus, \emph{From Prephysical Hypotheses to Physical Models: Collections} (81 pages) and \emph{Mathematical Supplement} (420 pages), supplied 10 November 2026. The scientific component titles and precise locations used in this assessment are identified in the text. These are same-program sources, not independent corroboration.
\bibitem{companion} K. Justus, \emph{From Prephysical Hypotheses to Physical Models: Mathematical Supplement. Calibration, identified sets and uniform limits} (10 November 2026), accompanying this assessment.
\end{thebibliography}
\clearpage
\phantomsection\addcontentsline{toc}{section}{Structure, observation and physical identification}
\includepdf[pages=1,pagecommand={\thispagestyle{empty}}]{../Structure_Observation_and_Physical_Identification.pdf}
\end{document}
